\documentclass[11pt,a4paper]{article}

\usepackage[margin=2.3cm]{geometry}
\usepackage{booktabs}
\usepackage{array}
\usepackage{hyperref}
\usepackage{xcolor}
\usepackage{longtable}
\usepackage{enumitem}
\usepackage{listings}

\hypersetup{
    colorlinks=true,
    linkcolor=black,
    urlcolor=blue
}

\lstset{
    basicstyle=\ttfamily\small,
    breaklines=true,
    frame=single
}

\title{\textbf{Computer Vision Project -- Exercise 04}\\Group Report}
\author{Group Submission}
\date{\today}

\begin{document}

\maketitle

\section{Scope of the Group Submission}

This report covers the group part of Exercise 04:

\begin{itemize}[noitemsep]
    \item Exercise 4.1: Face detection, tracking, and alignment
    \item Exercise 4.2: Face identification and open-set recognition
    \item Exercise 4.3: Face clustering and re-identification
    \item Exercise 4.4: DIR evaluation
\end{itemize}

Exercise 4.5 is treated as the individual challenge and is therefore not included in this group report.

\section{Face Detection, Tracking, and Alignment}

The face preprocessing pipeline consists of face detection, tracking, and alignment. MTCNN is used for face detection. The detected face is then aligned to a fixed resolution of $224 \times 224$ pixels before feature extraction.

For video processing, template matching was implemented in the \texttt{track\_face} method. The first frame is initialized using MTCNN. For later frames, a local search window around the previous face position is used. This avoids running MTCNN on every frame and makes the system closer to the intended video tracking pipeline.

\subsection{Tracking Drift and Fix}

During testing, a template drift problem was observed. The tracker sometimes slowly moved away from the actual face while still producing a high matching score. This especially affected the Alan\_Ball sequence, where recognition was initially very poor.

To reduce this problem, periodic MTCNN re-detection was added. The idea is to keep template tracking for efficiency, but re-anchor the tracker after a fixed number of frames.

The following re-detection intervals were tested:

\begin{table}[h]
\centering
\begin{tabular}{p{4cm}p{6cm}p{4cm}}
\toprule
\textbf{Setting} & \textbf{Description} & \textbf{Alan\_Ball Result} \\
\midrule
\texttt{redetect\_interval = 15} & Less frequent MTCNN re-detection & 277/432 \\
\texttt{redetect\_interval = 5} & Practical default balancing runtime and robustness & 385/432 \\
\texttt{redetect\_interval = 1} & MTCNN on every frame, highest accuracy but slowest & 426/432 \\
\bottomrule
\end{tabular}
\caption{Effect of periodic re-detection on the Alan\_Ball sequence.}
\end{table}

The final submitted group code uses:

\begin{lstlisting}
redetect_interval = 5
\end{lstlisting}

This was selected as the practical default because it still uses template tracking and avoids the high computational cost of running MTCNN on every frame. The \texttt{redetect\_interval = 1} result is reported as an additional accuracy-focused experiment.

\section{Face Identification}

Face identification was implemented using FaceNet embeddings and a k-nearest neighbor classifier. The implementation does not use sklearn for the group part.

The implemented components are:

\begin{itemize}[noitemsep]
    \item FaceNet embedding extraction from aligned faces
    \item Gallery construction using \texttt{partial\_fit}
    \item k-nearest neighbor prediction using Euclidean distances
    \item majority voting over the $k$ nearest neighbors
    \item posterior probability computed as $k_i/k$
    \item distance to the predicted class computed using the nearest neighbor from the predicted class
    \item color and grayscale embeddings stored for each training face
    \item prediction using both color and grayscale query embeddings
\end{itemize}

The final parameters are:

\begin{table}[h]
\centering
\begin{tabular}{lc}
\toprule
\textbf{Parameter} & \textbf{Value} \\
\midrule
\texttt{num\_neighbours} & 3 \\
\texttt{max\_distance} & 1.02 \\
\texttt{min\_prob} & 0.34 \\
\texttt{redetect\_interval} & 5 \\
\bottomrule
\end{tabular}
\caption{Final parameters used in the group source code.}
\end{table}

\section{Open-Set Recognition}

The open-set decision rule predicts a face as unknown if the distance to the predicted class is too high or the posterior probability is too low:

\begin{lstlisting}
if distance_to_prediction > max_distance or posterior < min_prob:
    prediction = "unknown"
\end{lstlisting}

The unknown identities used in testing were:

\begin{itemize}[noitemsep]
    \item Al\_Pacino
    \item Maria\_Shriver
\end{itemize}

A threshold sweep was performed after fixing the tracker. For \texttt{redetect\_interval = 5}, the best balanced threshold was:

\begin{table}[h]
\centering
\begin{tabular}{lccc}
\toprule
\textbf{Threshold} & \textbf{Known Accuracy} & \textbf{Unknown Rejection} & \textbf{Balanced Score} \\
\midrule
1.02 & 1040/1080 = 96.30\% & 154/154 = 100\% & 98.15\% \\
\bottomrule
\end{tabular}
\caption{Best threshold from the post-tracking-fix threshold sweep.}
\end{table}

\section{Identification Results}

The final practical code setting uses \texttt{redetect\_interval = 5}. Selected final-code results are:

\begin{table}[h]
\centering
\begin{tabular}{lc}
\toprule
\textbf{Identity} & \textbf{Result} \\
\midrule
Alan\_Ball & 385/432 = 89.12\% \\
Al\_Pacino & 48/49 rejected as unknown \\
Maria\_Shriver & 105/105 rejected as unknown \\
\bottomrule
\end{tabular}
\caption{Selected final-code identification results with \texttt{redetect\_interval = 5}.}
\end{table}

An additional high-accuracy experiment was performed with \texttt{redetect\_interval = 1}. This setting is slower because it effectively uses MTCNN every frame, but it shows the upper accuracy that can be reached when tracking drift is removed.

\begin{table}[h]
\centering
\begin{tabular}{lc}
\toprule
\textbf{Identity} & \textbf{Result} \\
\midrule
Alan\_Ball & 426/432 = 98.61\% \\
Manuel\_Pellegrini & 229/229 = 100\% \\
Marina\_Silva & 173/173 = 100\% \\
Nancy\_Sinatra & 67/73 = 91.78\% \\
Peter\_Gilmour & 173/173 = 100\% \\
Al\_Pacino & 49/49 rejected as unknown \\
Maria\_Shriver & 105/105 rejected as unknown \\
\bottomrule
\end{tabular}
\caption{Accuracy-focused experiment with \texttt{redetect\_interval = 1}.}
\end{table}

Overall result for the \texttt{redetect\_interval = 1} experiment:

\begin{itemize}[noitemsep]
    \item Known-person accuracy: 1068/1080 = 98.89\%
    \item Unknown rejection: 154/154 = 100\%
\end{itemize}

\section{Face Clustering and Re-Identification}

The \texttt{FaceClustering} class implements k-means clustering in FaceNet embedding space.

Implemented components:

\begin{itemize}[noitemsep]
    \item store unlabeled FaceNet embeddings
    \item estimate cluster centers with k-means
    \item store cluster assignments
    \item predict the closest cluster for a query face
    \item return the distance distribution to all cluster centers
\end{itemize}

Clustering was tested with two persons:

\begin{itemize}[noitemsep]
    \item Alan\_Ball
    \item Manuel\_Pellegrini
\end{itemize}

Saved clustering evidence includes:

\begin{itemize}[noitemsep]
    \item \texttt{results/models/clustering\_gallery\_final.pkl}
    \item \texttt{results/logs/cluster\_predictions.csv}
    \item \texttt{results/screenshots\_cluster/}
\end{itemize}

\section{DIR Evaluation}

The DIR evaluation was completed using the provided evaluation embeddings. The final metrics are:

\begin{table}[h]
\centering
\begin{tabular}{lc}
\toprule
\textbf{Metric} & \textbf{Value} \\
\midrule
DIR@FAR=1\% & 0.707091 \\
Threshold@1\% & -0.497051 \\
DIR@FAR=10\% & 0.883659 \\
Threshold@10\% & -0.646075 \\
Minimum FAR with DIR $\geq$ 90\% & 0.140605 \\
Threshold there & -0.671862 \\
\bottomrule
\end{tabular}
\caption{Final DIR evaluation results.}
\end{table}

The DIR curve was saved as:

\begin{lstlisting}
results/plots/dir_curve.png
\end{lstlisting}

\section{Conclusion}

The group part implements the full face recognition pipeline: detection, tracking, alignment, supervised face identification, open-set rejection, clustering-based re-identification, and DIR evaluation.

The most important improvement was fixing template-tracking drift through periodic MTCNN re-detection. The final source code uses \texttt{redetect\_interval = 5} as a practical runtime-aware setting, while the report also documents the higher-accuracy \texttt{redetect\_interval = 1} experiment.

\end{document}
